\chapter{极限：无限接近不是含糊其辞}

这一章要对付的，是数学里最会“耍赖”的一个词：无限。你早就见过它了——循环小数 $0.333\cdots$ 后面的省略号是无限，平行线“永不相交”是无限，连 $\sqrt{2}$ 的小数部分都拖着一条无限长的尾巴。可奇怪的是，初中数学一直在躲着它：老师会说“这个过程一直继续下去”，却很少说“继续下去到底得到了什么”。这一章我们不再躲。我们要看一看，两千多年前一个把人绕晕的悖论，是怎样逼着数学家发明出“极限”这个概念的；又是怎样把“无限接近”这句听起来模模糊糊的话，变成一条可以严格检查的规则。学完之后你会发现：无限并不神秘，神秘的只是没说清楚的无限。

\section{一场永远跑不完的比赛？}

想象一场特殊的赛跑。选手阿基里斯是传说中的飞毛腿，速度是每秒 $10$ 米；他的对手是一只乌龟，速度只有每秒 $1$ 米。为了公平起见（其实是故意找茬），乌龟先跑 $100$ 米，阿基里斯才从起点出发。

请问：阿基里斯能追上乌龟吗？

你的直觉一定会喊：当然能！他快十倍呢，追上不过是几秒钟的事。可两千多年前，哲学家芝诺却讲了一个让人心里发毛的推理。他说：好，我们一步一步来看。

第一步，阿基里斯先跑到乌龟出发时的位置，也就是跑完 $100$ 米。这需要 $10$ 秒。可就在这 $10$ 秒里，乌龟也没闲着，它往前爬了 $10$ 米，到达了 $110$ 米处。

第二步，阿基里斯再跑 $10$ 米，到达 $110$ 米处。这又花了 $1$ 秒。而乌龟又趁机爬了 $1$ 米，到了 $111$ 米处。

第三步，阿基里斯跑 $1$ 米，花 $0.1$ 秒，乌龟又爬了 $0.1$ 米……

发现规律了吗？每当阿基里斯到达乌龟“刚才所在的位置”，乌龟都已经又往前挪了一小段。这个过程可以一步接一步地说下去，第四步、第五步、第一万步，永远没有说完的时候。芝诺得意地宣布：既然追上的过程包含无限多个步骤，而每个步骤都要花一点时间，那么阿基里斯永远追不上乌龟。

这就是著名的芝诺悖论。说它“悖论”，不是因为它结论对，而是因为它看起来每一步都讲得通，结论却荒唐得离谱——谁都知道跑得快的人能追上爬得慢的乌龟。中国古代也有类似的说法，《庄子》里写着“一尺之棰，日取其半，万世不竭”：一根一尺长的木棒，每天取走它的一半，取一万年也取不完。东西方的古人不约而同地盯上了同一个怪物：无限分割。

\section{先把直觉押在桌上}

在学任何工具之前，先凭直觉猜一猜。下面三个说法，你觉得哪些对、哪些错？

第一，芝诺说的那无限多个时间段——$10$ 秒、$1$ 秒、$0.1$ 秒、$0.01$ 秒……加起来等于无穷大，所以追上要花无限长的时间。

第二，$0.999\cdots$（$9$ 无限循环）是一个比 $1$ 小一点点的数，它永远接近 $1$，但不等于 $1$。

第三，一根木棒每天取走一半，取无限多天之后，木棒会剩下“无穷小的一点点”，而不是正好取完。

先别急着翻页，把你的猜测记下来。这一章结束时，你要回来给它们改卷。

这三个问题其实问的是同一件事：无限多个“越来越小的量”堆在一起，结果到底是什么？是有限的，还是无限的？是“差一点点”，还是“正好相等”？靠拍脑袋解决不了，因为直觉在这里出了名的不可靠。我们需要一套能把“无限接近”说清楚的语言。

\section{第一次尝试：一项一项地加}

最朴素的办法是：别管“无限”，我们就老老实实地加。芝诺那堆时间段是
\[
10+1+0.1+0.01+0.001+\cdots
\]
我们把省略号砍掉，只加前面几项，看看会得到什么：
\[
\begin{aligned}
10 &= 10,\\
10+1 &= 11,\\
10+1+0.1 &= 11.1,\\
10+1+0.1+0.01 &= 11.11,\\
10+1+0.1+0.01+0.001 &= 11.111.
\end{aligned}
\]
有意思。每多加一项，结果就多在末尾添一个 $1$，一路朝着 $11.111\cdots$ 这个方向爬过去。把这张“爬坡记录”整理成表格，趋势就更清楚了：

\begin{center}
\begin{tabular}{ccc}
\toprule
加的项数 $n$ & 部分和 $S_n$ & 与 $\dfrac{100}{9}$ 的差距 \\
\midrule
$1$ & $10$ & $1.111\cdots$ \\
$2$ & $11$ & $0.111\cdots$ \\
$4$ & $11.11$ & $0.00111\cdots$ \\
$8$ & $11.111111$ & 约 $1.1\times 10^{-7}$ \\
\bottomrule
\end{tabular}
\end{center}

每多加两项，差距的小数点后就多出一个 $0$。而且 $11.111\cdots$ 这个无限循环小数，恰好等于分数 $\dfrac{100}{9}$，约等于 $11.11$ 秒。这不正是“追上所需的时间”该有的答案吗？用路程除以相对速度也能验证：阿基里斯每秒比乌龟多跑 $9$ 米，要追回 $100$ 米的差距，需要 $\dfrac{100}{9}$ 秒。两条路指向同一个数，这让人松了一口气。

但先别庆祝，这个办法有两个大窟窿。

第一个窟窿：我们永远加不完。“只加前几项”得到的 $11.11$、$11.111$，每一个都只是近似值。凭什么说无限多项加起来的结果，就正好等于 $\dfrac{100}{9}$，而不是比它大一点点或小一点点？我们刚才实际上是“猜”出了答案，然后用另一种方法（相对速度）验证它猜得对。可如果换一道题，没有别的办法可以验证呢？猜，终究不是数学。

第二个窟窿更隐蔽：我们偷偷用了一个没证明的信念——“结果一路朝着 $\dfrac{100}{9}$ 爬，最后就到 $\dfrac{100}{9}$”。“一路朝着”和“最终到达”之间，隔着整整一个无限。直觉觉得这条缝可以跳过去，但数学要求我们把桥架起来。

\begin{fanli}
警惕这样一种说法：“反正加到后面每一项都小得可以忽略，所以总和就是前面那些项的值，后面不算也罢。”这在小数计算里或许够用，但它恰恰把真正的问题藏起来了。“小得可以忽略”到底是多小？是对谁来说可以忽略？如果你只关心秒的精度，$0.001$ 秒可以忽略；如果你在调度一枚火箭，$0.001$ 秒可能意味着几十米的偏差。“忽略”是一个工程决定，不是一个数学结论。数学要回答的是：把无限多项一项不丢地全算进去，结果到底存不存在，等于多少。
\end{fanli}

\section{给“无限接近”下定义}

要架桥，先要有桥墩。第一个桥墩是把“一串无穷无尽、但有先后顺序的数”当作一个整体来研究。

\begin{dingyi}[数列]
按照某种规律排成一队、可以一直写下去的数，叫做一个\textbf{数列}。排在第 $n$ 个位置的数叫做数列的第 $n$ 项，记作 $a_n$。整个数列可以写成
\[
a_1,\ a_2,\ a_3,\ \cdots,\ a_n,\ \cdots
\]
\end{dingyi}

比如我们刚才加出来的那些“部分和”，就是一个数列：
\[
10,\ 11,\ 11.1,\ 11.11,\ 11.111,\ \cdots
\]
它的第 $n$ 项是“前 $n$ 个时间段之和”。注意数列和集合不一样：数列里的数有先后，可以重复出现，而且永远有“下一项”。研究无限过程的办法，就是把过程切成一站一站的，把每一站的结果排成数列，然后问：这个数列的“终点”是什么？

可是“终点”这个词本身就需要澄清。数列没有最后一项，哪来的终点？我们需要换一种问法，这就是极限思想的关键一跃：\textbf{不问数列“最后到达哪里”，而问数列“无限接近哪个数”}。那什么叫“无限接近”？“接近”是距离小，“无限接近”难道是“距离无限小”？“无限小”又是什么？我们刚说过，不能再用模糊的词去定义模糊的词。

数学家想出的办法非常聪明，他们把“无限接近”换成了一句可以逐字检查的话。先给你一个日常版本：

\begin{quote}
你想要多近，我就能让你看到多近——而且从那一步往后，永远保持那么近。
\end{quote}

把它写成数学语言，就是数列极限的定义。

\begin{dingyi}[数列的极限，直观版本]
设有一个数列 $a_1,a_2,a_3,\cdots$ 和一个数 $L$。如果无论你指定一个多么小的正数 $d$，我都能找出数列中的某一个位置 $N$，使得从第 $N$ 项往后的\textbf{所有}项，与 $L$ 的距离都小于 $d$，也就是
\[
|a_n-L|<d \quad (n>N),
\]
那么就说这个数列\textbf{收敛}到 $L$，$L$ 叫做这个数列的\textbf{极限}，记作
\[
\lim_{n\to\infty}a_n=L.
\]
\end{dingyi}

这个定义值得逐句拆开看，因为它是全书的重点之一。

首先，它是一场你出题、我答题的游戏。你负责挑一个“接近标准”$d$——可以是 $0.1$，可以是 $0.0001$，可以是亿分之一，随便你挑多小。我负责证明：数列从某一项开始，所有项都落在以 $L$ 为中心、宽度为 $2d$ 的“走廊”$(L-d,\,L+d)$ 里，而且再也不跑出去。如果你挑一个更小的 $d$，我可能要把位置 $N$ 往后挪——比如从第 $10$ 项挪到第 $10000$ 项——但只要你肯出题，我就总能答出来。只要这场游戏我永远赢，极限就是 $L$。

其次，注意定义里没有一个字提到“无限小”，也没有提到“无限接近”。它只用到了普通的不等式 $|a_n-L|<d$，而 $d$ 是一个普普通通的、大于零的数。妙处全在“无论多么小”这五个字里：它用一个一个有限的、可以检查的挑战，代替了对“无限”的直接谈论。这就是把含糊的“无限接近”说清楚的办法——不是给你一个神秘的无穷小量，而是给你一个永不失败的应答程序。

这套语言来得并不容易。微积分刚诞生的时候，数学家们直接使用“无穷小量”进行计算，得出的结果往往正确，却没人说得清无穷小到底是什么。批评者讥讽它是“消失了的量的鬼魂”：一个量一会儿不是零（否则不能当除数），一会儿又是零（否则舍不掉），自相矛盾。争论持续了上百年，直到十九世纪，数学家们才下定决心把“无穷小”这个说不清的角色请出舞台，换上刚才那种用不等式说话的走廊游戏。从那以后，“无限接近”不再是一种感觉，而是一条可以逐字检查的规则。你在这一章见到的定义，正是那场漫长争论的终点。

\begin{shuyu}{极限}
\textbf{直观解释}：数列的极限就是它“想去的那个地方”。不需要真的到达，但从某一站开始，它会一直待在那个地方的任意小的附近。\\
\textbf{正式定义}：对任意正数 $d$，都存在位置 $N$，使得 $n>N$ 时 $|a_n-L|<d$。\\
\textbf{一个例子}：数列 $1,\dfrac{1}{2},\dfrac{1}{3},\dfrac{1}{4},\cdots$ 收敛到 $0$。你出 $d=0.01$，我答 $N=100$，因为从第 $101$ 项起每一项都小于 $0.01$；你出更小的 $d$，我取一个比 $\dfrac{1}{d}$ 大的 $N$ 就行。
\end{shuyu}

用数轴画出来更清楚。下面这条数轴上，黑点依次是数列 $\dfrac{1}{2},\dfrac{3}{4},\dfrac{7}{8},\dfrac{15}{16},\cdots$ 的前几项。它们挤向 $1$，而以 $1$ 为中心的走廊（虚线之间）越来越窄，但只要走廊还有宽度，就总有某一项之后全体“入住”。

\begin{center}
\begin{tikzpicture}[scale=1]
  \draw[->] (-0.4,0) -- (7.6,0);
  \foreach \x in {0,3.5,5.25,6.125,6.5625,7}
    \draw (\x,0.08) -- (\x,-0.08);
  \node[below] at (0,-0.1) {$0$};
  \node[below] at (3.5,-0.1) {$\frac{1}{2}$};
  \node[below] at (5.25,-0.1) {$\frac{3}{4}$};
  \node[below] at (6.65,-0.45) {$\frac{7}{8},\ \frac{15}{16}\cdots$};
  \node[below] at (7,-0.1) {$1$};
  \fill (3.5,0) circle (2pt);
  \fill (5.25,0) circle (2pt);
  \fill (6.125,0) circle (2pt);
  \fill (6.5625,0) circle (2pt);
  \fill (6.78,0) circle (1.6pt);
  \fill (6.89,0) circle (1.2pt);
  \draw[dashed] (6.3,0.5) -- (6.3,-0.25);
  \draw[dashed] (7,0.55) -- (7,0.5);
  \node[above] at (6.65,0.55) {以 $1$ 为中心的走廊};
\end{tikzpicture}
\end{center}

还要注意一个容易误会的地方：收敛不等于到达。数列 $\dfrac{1}{n}$ 的每一项都大于 $0$，没有任何一项等于 $0$，但它的极限就是 $0$。极限描述的是趋势和归宿，不要求“撞线”。反过来，有些数列压根没有归宿：$1,-1,1,-1,\cdots$ 在两个值之间反复横跳，你随便画一条宽度小于 $2$ 的走廊，无论套在 $1$ 上还是套在 $-1$ 上，数列都会不断跑出走廊。这样的数列叫做\textbf{发散}，它没有极限。

\begin{fanli}
再拆一个更精致的骗局。有人说：“数列 $a_n$ 离 $L$ 的距离一项比一项小，所以 $L$ 就是它的极限。”对吗？看反例：数列 $\dfrac{1}{n}$（也就是 $1,\frac12,\frac13,\cdots$）到 $-1$ 的距离是 $1+\dfrac{1}{n}$，确实一项比一项小！按那个说法，$-1$ 应该是它的极限。可它的极限明明是 $0$。错在哪？错在“距离越来越小”太弱了——$1+\dfrac{1}{n}$ 虽然一直在缩小，却始终大于 $1$，根本缩不到“任意小”。极限定义里“无论 $d$ 多小都能做到”这一关，它过不去。这个反例说明：\textbf{单调地靠近一个目标，和任意地接近一个目标，是两回事}。定义里每一个字都不是摆设。
\end{fanli}

\section{核心定理：无限个越来越小的数可以加出确定的和}

有了极限的语言，我们终于可以正面回答芝诺的问题了。先提出一个准确的命题。

\begin{dingli}[无穷等比数列的和]\label{ch9:geo}
设首项为 $a$、公比为 $r$ 的等比数列 $a,ar,ar^2,ar^3,\cdots$，其中 $|r|<1$。把前 $n$ 项加起来得到部分和
\[
S_n=a+ar+ar^2+\cdots+ar^{n-1}.
\]
那么当 $n$ 无限增大时，$S_n$ 收敛，且
\[
\lim_{n\to\infty}S_n=\frac{a}{1-r}.
\]
这个极限就叫做这个无穷等比数列各项的和，记作
\[
a+ar+ar^2+\cdots=\frac{a}{1-r}.
\]
\end{dingli}

注意这个定理的措辞：它没有说“把无限多项加起来”——那种操作我们根本不会做。它说的是：有限项的和 $S_n$ 组成一个数列，这个数列收敛，我们把它的极限\textbf{定义}为无穷和。无限求和不是我们真的会做的加法，而是极限的别名。这个区别正是解开芝诺悖论的钥匙。

\begin{proof}
证明分两步。第一步，求出 $S_n$ 的紧凑表达式，这是初中代数里的老技巧“错位相减”。由
\[
S_n=a+ar+ar^2+\cdots+ar^{n-1},
\]
两边同乘 $r$，得
\[
rS_n=ar+ar^2+\cdots+ar^{n-1}+ar^{n}.
\]
两式相减，中间所有项全部抵消，只剩首尾：
\[
S_n-rS_n=a-ar^n,
\qquad\text{即}\qquad
S_n=\frac{a(1-r^n)}{1-r}.
\]
第二步，用极限的定义看 $n$ 增大时会发生什么。把上式改写成
\[
S_n=\frac{a}{1-r}-\frac{a}{1-r}\,r^n.
\]
因为 $|r|<1$，$r^n$ 的绝对值随着 $n$ 增大而不断缩小：比如 $r=0.1$ 时，$r^n$ 依次是 $0.1,0.01,0.001,\cdots$。任给一个正数 $d$，只要 $n$ 足够大，$|r^n|$ 就能小到让 $\left|\dfrac{a}{1-r}\right|\cdot|r^n|$ 小于 $d$。也就是说 $S_n$ 与 $\dfrac{a}{1-r}$ 的距离可以任意小，按定义
\[
\lim_{n\to\infty}S_n=\frac{a}{1-r}.
\]
两步都是普通的代数与不等式推理，合在一起却跨越了无限。
\end{proof}

公式里的条件 $|r|<1$ 不是摆设。如果 $r=2$，数列的和是 $1+2+4+8+\cdots$，越加越大，部分和一路奔向无穷，根本不收敛；如果 $r=-1$，和式变成 $1-1+1-1+\cdots$，部分和在 $1$ 和 $0$ 之间来回跳，同样没有极限。这两种情况都发散，$\dfrac{a}{1-r}$ 这个公式一概不能套用。数学定理的条件就像电器的额定电压：插错了，再漂亮的公式也会烧掉。

这个证明里有一个值得停顿的地方：第二步为什么可信？因为它每一步都可以拿出来检查。你不信 $r^n$ 能任意小？取 $r=\dfrac{1}{2}$，你出 $d=10^{-9}$，我只要取 $n=31$，就有 $2^{-31}<10^{-9}$。你出更苛刻的 $d$，我给出更大的 $n$。没有玄学，只有应答。

\begin{tuilun}[芝诺悖论的解决]
阿基里斯追上乌龟所需的时间，是各阶段时间之和的极限：
\[
10+1+0.1+0.01+\cdots=\frac{10}{1-0.1}=\frac{100}{9}\ \text{秒}.
\]
无限多个阶段不意味着无限长时间。芝诺推理的真正错误在于他默认了“无限多个正数相加必得无穷大”，而刚才的定理给出了反例。
\end{tuilun}

还有一个更漂亮的看法：不套公式，直接看图。把一个面积为 $1$ 的正方形，先取走一半，再把剩下部分取走一半，再取一半……每次取走的面积依次是 $\dfrac{1}{2},\dfrac{1}{4},\dfrac{1}{8},\dfrac{1}{16},\cdots$。从图里你能直接“看见”这些块正好拼满整个正方形：

\begin{center}
\begin{tikzpicture}[scale=5]
  \draw (0,0) rectangle (1,1);
  \draw (0.5,0) -- (0.5,1);
  \draw (0.5,0.5) -- (1,0.5);
  \draw (0.75,0) -- (0.75,0.5);
  \draw (0.75,0.25) -- (1,0.25);
  \node at (0.25,0.5) {$\frac{1}{2}$};
  \node at (0.75,0.75) {$\frac{1}{4}$};
  \node at (0.625,0.25) {$\frac{1}{8}$};
  \node at (0.875,0.375) {$\frac{1}{16}$};
  \node at (0.875,0.1) {$\cdots$};
\end{tikzpicture}
\end{center}

图直观地说：$\dfrac{1}{2}+\dfrac{1}{4}+\dfrac{1}{8}+\cdots=1$，和公式 $\dfrac{a}{1-r}$ 取 $a=\dfrac{1}{2}$、$r=\dfrac{1}{2}$ 的结果完全一致。这也顺手回答了开头的第三个问题：“日取其半”的木棒，每天取下的长度加起来恰好等于 $1$ 尺——无限天之后取走的总量正好是一整根，没有“剩下一点点”这回事。至于“每天取走一半，永远还有一半剩着”，那是过程永远不完结的意思；而“总量恰好是 $1$”说的是极限。两句话不矛盾，因为它们谈论的是不同的对象：一个是过程，一个是极限。把这两个对象混为一谈，正是芝诺悖论的病灶。

\section{换一个视角：数轴上有没有洞}

到目前为止我们一直在问“极限是多少”。现在换一个角度，问一个更深的问题：极限一定存在吗？它落在哪种数里？

看一个例子。数列
\[
1,\ 1.4,\ 1.41,\ 1.414,\ 1.4142,\ \cdots
\]
每一项都是有限小数，也就是有理数。它们越挤越拢：从某一项起，任意两项之间的距离可以任意小（因为它们的小数位相同的部分越来越长）。按直觉，它们应该挤向数轴上的某个点。可是我们都知道，这个点对应的长度——边长为 $1$ 的正方形的对角线——是 $\sqrt{2}$，而 $\sqrt{2}$ 不是有理数。如果我们只允许有理数存在，那么这个数列拼命地想收敛，却找不到归宿：它的极限“不在城里”。

这就像一部电梯匀速上升，却一直到达不了任何楼层——问题不在电梯，在大楼少盖了一层。有理数这根数轴上布满了这样的“洞”：随便用圆规在数轴上截一段对角线，落点就掉进一个洞里。

\begin{shuyu}{实数的完备性}
\textbf{直观解释}：实数轴上没有洞。任何一个“越挤越拢”的数列（后文叫它柯西数列），都能找到一个实数作为它的极限。\\
\textbf{正式定义}：实数系满足完备性公理——每个柯西数列都收敛到某个实数；等价地，每个有上界的实数集合都有最小上界。\\
\textbf{一个例子}：$\sqrt{2}$ 就是数列 $1,1.4,1.41,\cdots$ 的极限；不是先知道 $\sqrt{2}$ 再找数列，而是可以反过来，用这样的数列把 $\sqrt{2}$ “定义”出来。
\end{shuyu}

完备性是一条公理——不是我们证明出来的，而是我们要求实数满足的性质。你可以把它理解成盖楼时的承诺：只要数列“应该”有极限，我们就保证数轴上有这个点，绝不缺货。正是在这个承诺之上，极限理论才站得稳，微积分的大厦才盖得起来。

完备性还有一个更朴素的化身，值得单独看一眼。想象一个贴着天花板不断吹大的气球：它每一步都在变大，却永远越不过天花板。直觉告诉我们，它只能越来越紧地贴向某个确定的高度。换成数列的语言就是：\textbf{一个每一项都不小于前一项、又全体不超过某个固定数的数列，必定收敛}。这叫做单调有界收敛原理，它和上面的柯西准则一样，都是“数轴没有洞”的不同说法。后文实验三里那个 $\left(1+\dfrac{1}{n}\right)^n$，正是靠这条原理安身立命的：可以证明它单调递增，并且永远小于 $3$，所以它的极限必然存在——数学家是先证明了极限存在，才给这个极限起名叫 $e$ 的。先有存在性，后有名字，这个顺序值得记住。

现在回头处理开头第二个问题：$0.999\cdots$ 到底等不等于 $1$？先用定理 \ref{ch9:geo} 算一算。$0.999\cdots$ 的意思是无穷和
\[
\frac{9}{10}+\frac{9}{100}+\frac{9}{1000}+\cdots,
\]
首项 $a=\dfrac{9}{10}$，公比 $r=\dfrac{1}{10}$，所以
\[
0.999\cdots=\frac{\frac{9}{10}}{1-\frac{1}{10}}=\frac{\frac{9}{10}}{\frac{9}{10}}=1.
\]
没有“差一点点”，就是 $1$。觉得别扭吗？别扭的根源在于把 $0.999\cdots$ 想成“一个正在往 $1$ 爬、永远爬不到的数”。但它不是过程，它是那个无穷和的极限值，而极限值就是 $1$。换句话说：$1$ 和 $0.999\cdots$ 是同一个实数的两种写法，就像 $\dfrac{1}{2}$ 和 $\dfrac{2}{4}$ 是同一个分数的两种写法一样。事实上每一个有限小数都有这样一对“双胞胎写法”，比如 $0.5=0.4999\cdots$。

\begin{lianjie}
极限的想法并不只属于纯数学。你在科学课上学过的“理想气体”“光滑斜面”“瞬时速度”，全是极限的产物：真实气体分子有大小，理想气体是把分子大小趋向 $0$ 的极限；真实斜面有摩擦，光滑斜面是摩擦趋向 $0$ 的极限。工程师在计算机里模拟天气、设计芯片，也靠极限的反面操作——把连续变化切成有限多步来算，步数越多越接近真相。你手机里的每一场天气预报，都是“有限步逼近无限”的产物。这一章学的极限，正是下一章导数的地基：导数就是“平均变化率当时间间隔趋向 $0$ 时的极限”。
\end{lianjie}

\begin{toolbox}
本章新添的工具：
\begin{itemize}
  \item \textbf{数列}：把无限过程切成一站一站，研究这个“站队”而不是空想无限本身。
  \item \textbf{极限的定义（走廊游戏）}：对任意正数 $d$，存在位置 $N$，$n>N$ 时 $|a_n-L|<d$。用它可以把“无限接近”翻译成可检查的不等式。
  \item \textbf{无穷等比数列求和}：$|r|<1$ 时 $a+ar+ar^2+\cdots=\dfrac{a}{1-r}$。无穷和是部分和的极限，不是“真的加无限多次”。
  \item \textbf{完备性的观念}：实数轴没有洞，该收敛的数列都有归宿。
\end{itemize}
\end{toolbox}

\section{数学实验}

\begin{shiyan}
\textbf{实验一：亲手“加”出一个无穷和。}找一张计算器，依次计算
\[
\frac{1}{2},\quad \frac{1}{2}+\frac{1}{4},\quad \frac{1}{2}+\frac{1}{4}+\frac{1}{8},\quad \cdots
\]
每加一项记录一次结果，至少加到 $\dfrac{1}{2^{10}}=\dfrac{1}{1024}$。你会得到 $0.5,0.75,0.875,0.9375,\cdots$。观察：每加一项，结果与 $1$ 的距离缩小为原来的一半。再算一算，加到第几项时结果与 $1$ 的距离小于 $0.001$？对照走廊游戏，你出的 $d=0.001$ 对应的 $N$ 是多少？

\textbf{实验二：给 $\sqrt{2}$ 修一条“逼近公路”。}不用开方键，只用四则运算。从区间 $[1,2]$ 开始，计算中点 $1.5$，因为 $1.5^2=2.25>2$，所以 $\sqrt{2}$ 在 $[1,1.5]$ 里；再取中点 $1.25$，平方后与 $2$ 比较，决定取左半还是右半……每做一次，区间长度减半。做 $8$ 次，记下每次的区间端点，你会得到 $\sqrt{2}\approx 1.414$。想想：这个“区间套”过程和 $\sqrt{2}$ 的小数展开 $1,1.4,1.41,\cdots$ 有什么共同点？如果数轴上有洞，这个实验还能保证夹出一个数吗？

\textbf{实验三（选做）：一个会变魔术的数。}计算 $\left(1+\dfrac{1}{n}\right)^n$：$n=1$ 时是 $2$，$n=10$ 时约 $2.594$，$n=100$ 时约 $2.705$，$n=1000$ 时约 $2.717$。它似乎在收敛，而且收敛到一个以 $2.718$ 开头的神秘数。先记下这个观察，本章末尾的思考题会再回到它。
\end{shiyan}

\section*{思考题}

\textbf{观察题}
\begin{sikao}
  \item 观察数列 $\dfrac{1}{2},\dfrac{2}{3},\dfrac{3}{4},\dfrac{4}{5},\cdots$（第 $n$ 项是 $\dfrac{n}{n+1}$）。先算前几项猜出它的极限，再算一算第 $n$ 项与这个极限的距离是多少。\textbf{提示}：把 $\dfrac{n}{n+1}$ 改写成 $1-\dfrac{1}{n+1}$，距离就一目了然了。
  \item 用本章定理把 $0.333\cdots$ 和 $0.121212\cdots$ 分别写成分数。观察第二个数：它的“公比”是多少？\textbf{提示}：$0.121212\cdots$ 每两位循环一次，可以看成首项 $\dfrac{12}{100}$、公比 $\dfrac{1}{100}$ 的无穷等比数列。
\end{sikao}

\textbf{证明题}
\begin{sikao}
  \item 用走廊游戏严格证明：数列 $\dfrac{1}{n}$ 收敛到 $0$。也就是对任意正数 $d$，明确写出你要取的 $N$（可以依赖于 $d$），并说明为什么 $n>N$ 时 $\left|\dfrac{1}{n}-0\right|<d$。\textbf{提示}：你需要的只是 $\dfrac{1}{n}<d$，两边取倒数看看 $n$ 该多大。
  \item 证明一个数列不可能同时收敛到两个不同的数 $L$ 和 $M$。\textbf{提示}：设 $L$ 和 $M$ 之间的距离为 $h>0$。取走廊宽度 $d=\dfrac{h}{2}$，分别以 $L$、$M$ 为中心画走廊——这两条走廊互不重叠。如果数列最终同时待在两条走廊里，会发生什么矛盾？
\end{sikao}

\textbf{挑战题}
\begin{sikao}
  \item 实验三里那个数列 $\left(1+\dfrac{1}{n}\right)^n$ 看起来在收敛。查一查或者想一想：它的极限叫做自然常数 $e$，约等于 $2.71828$。有趣的是，这个数列每一项都是有理数，极限 $e$ 却是无理数。结合本章关于完备性的讨论，说说这件事为什么一点也不奇怪。\textbf{提示}：回忆 $\sqrt{2}$ 的例子——“每一项是哪种数”和“极限是哪种数”是两个独立的问题。
  \item 芝诺还有第二个悖论“飞矢不动”：飞行的箭在每一瞬间都占据一个确定的位置，所以它在每一瞬间都是静止的，因此箭从未运动。用本章的语言指出这个推理的漏洞。\textbf{提示}：想想“位置在每一瞬间确定”和“没有运动”之间的跳跃，是不是又把“过程”和“极限”搞混了？“某一瞬间的速度”到底该如何定义？这正是下一章要解决的问题。
\end{sikao}

\begin{pandeng}
继续往上爬，你会遇到这些风景：把本章的 $d$ 换成希腊字母 $\varepsilon$、把 $N$ 的选取写成严格证明，就是大学数学分析里的“$\varepsilon$--$N$ 语言”；“越挤越拢的数列必有极限”的严格表述叫柯西收敛准则，它和“单调有界数列必收敛”等价，都是实数完备性的化身；无穷等比数列推广成一般的无穷级数，判别它们收敛还是发散是级数理论的开端；而 $\left(1+\frac{1}{n}\right)^n$ 的极限 $e$ 会在对数、复利和微分方程里反复登场。竞赛中常见的“放缩法”求极限、以及“夹逼定理”（两边都挤向 $L$，中间无处可逃），本质上都是走廊游戏的变体。
\end{pandeng}

\section*{通往下一章的门}

这一章我们驯服了“无限多个量的累加”。但芝诺的飞矢还悬在半空，留下一个更刁钻的问题：什么叫“某一瞬间”的速度？

汽车仪表盘显示时速 $60$ 公里，这个数是哪儿来的？速度是路程除以时间——可“一瞬间”的时间里，汽车走的路程是 $0$，时间也是 $0$，$\dfrac{0}{0}$ 是什么？都不是。直觉又失灵了。

出路你已经见过雏形了：别问“那一瞬间”，改问“从那时刻起的一小段时间”。算一算 $1$ 秒内的平均速度、$0.1$ 秒内的、$0.01$ 秒内的……如果这些平均速度排成的数列收敛，它的极限就有资格被称作“瞬时速度”。同样地，曲线上某一点的切线斜率，是割线斜率当两点靠拢时的极限。

平均变化率的极限，就是导数。下一章，我们把走廊游戏搬进函数的世界，去看看“变化”本身如何被精确地度量。
